\documentclass[border=5mm]{standalone}
% ==============================================================================
% SETUP.TEX - CONFIGURACIÓN GLOBAL DEL PROYECTO
% ==============================================================================
% Este archivo centraliza todos los paquetes y estilos.
% Debe incluirse en el preámbulo del libro principal y de los archivos standalone.
% \PassOptionsToPackage{gray}{xcolor}
% --- 1. CODIFICACIÓN E IDIOMA ---
% fix-cm habilita las fuentes Computer Modern en tamaños arbitrarios (por debajo
% de 5 pt, entre otros). Debe cargarse antes que fontenc.
\usepackage{fix-cm}
\usepackage[utf8]{inputenc}
\usepackage[T1]{fontenc}
\usepackage[spanish, es-tabla]{babel}  
\usepackage{csquotes} % Recomendado para babel y citas

\usepackage{import}
% mode=image: Usa el PDF pre-compilado, nunca intenta recompilar.
% Debes compilar cada figura manualmente antes de compilar el libro.
\usepackage[mode=image, subpreambles=false]{standalone}

% --- 2. MATEMÁTICAS Y CIENCIAS ---
\usepackage{amsmath, amssymb, amsfonts, mathtools}
\usepackage{enumitem}

% LaTeX trae \DeclareMathSizes{5}{5}{5}{5}, así que a 5 pt (\tiny) los subíndices
% salen del mismo tamaño que la base: en las etiquetas de figuras el M de
% $\omega_M$ queda desproporcionado. Se restablece la proporción habitual.
\DeclareMathSizes{5}{5}{3.7}{3}

% --- 3. DISEÑO Y GEOMETRÍA --- 
% Unificamos los márgenes para todo el documento
\usepackage[left=2.5cm,right=2.5cm,top=3cm,bottom=3cm]{geometry}
\makeatletter
\ifstandalone % Evita que geometry dañe el recorte de las figuras bajándolas 3cm
  \@ifclasswith{standalone}{crop=false}{
    % Es un capítulo/sección: Dejamos que geometry actúe.
  }{
    % Es una figura: Neutralizamos geometry para que no las recorte mal.
    \geometry{pass}
  }
\fi
\makeatother
\usepackage{parskip} % Espaciado entre párrafos sin sangría
\usepackage{float}   % Para usar [H] en figura

% Ajustes de flotantes para evitar espacios verticales exagerados
\renewcommand{\topfraction}{0.95}      % Permite que las figuras ocupen hasta el 95% superior
\renewcommand{\bottomfraction}{0.95}   % Permite que las figuras ocupen hasta el 95% inferior
\renewcommand{\textfraction}{0.05}     % Permite hojas con tan solo un 5% de texto
\renewcommand{\floatpagefraction}{0.8} % Requiere que una página de solo figuras esté 80% llena
\setcounter{topnumber}{4}
\setcounter{bottomnumber}{4}
\setcounter{totalnumber}{10}

% --- 4. GRÁFICOS Y TABLAS ---
\usepackage{graphicx}
\usepackage{booktabs} % Tablas profesionales
\usepackage{caption}  % Control de captions y \captionof
\usepackage{tikz}
\usetikzlibrary{arrows.meta, calc, angles, quotes, babel, backgrounds}
\usepackage{pgfplots}
\usepgfplotslibrary{groupplots}
\usepgfplotslibrary{external}
\usepgfplotslibrary{patchplots}
\pgfplotsset{compat=1.18}
\usepackage[american, straightvoltages]{circuitikz}

% --- 5. COLORES Y CAJAS ---

\usepackage[table]{xcolor}
\usepackage{tcolorbox}

% Definición de colores corporativos/estilo
\definecolor{utnblue}{RGB}{0, 51, 102}
\definecolor{codegreen}{rgb}{0,0.6,0}
\definecolor{codegray}{rgb}{0.5,0.5,0.5}
\definecolor{codepurple}{rgb}{0.58,0,0.82}
\definecolor{backcolour}{rgb}{0.96,0.96,0.96}

% --- 6. ENTORNOS PERSONALIZADOS (CAJAS) ---

% Caja "Resultado" (Azul Corporativo) — Definiciones, fórmulas clave, resultados importantes.
% Uso: \begin{resultado}[Fórmula de Euler] ... \end{resultado}
\newtcolorbox{resultado}[1][]{
    colback=utnblue!5!white,
    colframe=utnblue,
    title=\textbf{#1},
    fonttitle=\bfseries,
    sharp corners,
    boxrule=0.5mm
}

% Caja "Nota" (Gris) — Advertencias, aplicaciones, contexto secundario.
% Uso: \begin{nota}[Cuidado con la calculadora] ... \end{nota}
% Uso: \begin{nota} ... \end{nota}  → título por defecto "Nota"
\newtcolorbox{nota}[1][Nota]{
    colback=codegray!5!white,
    colframe=codegray,
    title=\textbf{#1},
    fonttitle=\bfseries,
    sharp corners,
    boxrule=0.5mm
}

% Caja inline para resaltar un valor y enlazarlo con su otra aparición en una tabla
% o en una cuenta: mismo color, mismo valor.
% Uso: \marcavalor{$4$}  o  \marcavalor[codegreen]{$4$}
\newtcbox{\marcavalor}[1][utnblue]{
    on line,
    boxsep=1pt, left=2pt, right=2pt, top=1pt, bottom=1pt,
    colback=#1!10!white,
    colframe=#1,
    boxrule=0.4pt,
    arc=2pt
}

% --- 7. CÓDIGO FUENTE (LISTINGS) ---
\usepackage{listings}

\lstdefinestyle{miestilo}{
    backgroundcolor=\color{backcolour},   
    commentstyle=\color{codegreen},
    keywordstyle=\color{magenta},
    numberstyle=\tiny\color{codegray},
    stringstyle=\color{codepurple},
    basicstyle=\ttfamily\footnotesize,
    breakatwhitespace=false,         
    breaklines=true,                 
    captionpos=b,                    
    keepspaces=true,                 
    numbers=left,                    
    numbersep=5pt,                  
    showspaces=false,                
    showstringspaces=false,
    showtabs=false,                  
    tabsize=2,
    frame=single,                    
    rulecolor=\color{codegray}
}

\lstset{style=miestilo}
% Soporte para tildes y ñ en bloques de código
\lstset{literate={ñ}{{\~n}}1 {á}{{\'a}}1 {é}{{\'e}}1 {í}{{\'i}}1 {ó}{{\'o}}1 {ú}{{\'u}}1}

% Operadores matemáticos personalizados

% Vector en sentido discreto (lista finita de coordenadas): negrita + flecha.
% Uso: \vect{v}, \vect{e}_k, \vect{0}.
\newcommand{\vect}[1]{\vec{\mathbf{#1}}}

% Fecha de versión y comando para capítulos standalone
\providecommand{\verdate}{\the\year.\ifnum\month<10 0\fi\the\month.\ifnum\day<10 0\fi\the\day}
\newcommand{\chapterversion}{%
    \vspace{-1.5cm}\begin{flushright}\small\textit{\color{codegray}Versión~\verdate}\end{flushright}\vspace{0.5cm}%
}

% --- 8. HIPERVÍNCULOS (DEBE SER EL ÚLTIMO PAQUETE) ---
\usepackage[colorlinks=true, linkcolor=utnblue, urlcolor=utnblue, citecolor=utnblue]{hyperref}

% --- 9. REFERENCIAS CRUZADAS ENTRE CAPÍTULOS STANDALONE ---
% Cuando se compila un capítulo standalone, xr-hyper lee los .aux de los
% otros capítulos (si existen) para resolver \ref{} a labels externos.
% En la compilación del libro completo este bloque no se activa.
\ifstandalone
  \usepackage{xr-hyper}
  \makeatletter
  \newcommand{\xrchap}[1]{%
    \edef\xrc@self{\detokenize{Capitulo#1}}%
    \edef\xrc@job{\jobname}%
    \ifx\xrc@self\xrc@job\else
      \IfFileExists{../#1capitulo/Capitulo#1.aux}%
        {\externaldocument{../#1capitulo/Capitulo#1}}{}%
    \fi
  }
  \makeatother
  \xrchap{01}\xrchap{02}\xrchap{03}\xrchap{04}
  \xrchap{05}\xrchap{06}\xrchap{07}\xrchap{08}
  \xrchap{09}\xrchap{10}\xrchap{11}\xrchap{12}
  \xrchap{13}
\fi
\usetikzlibrary{shadows.blur}

\begin{document}
    \begin{tikzpicture}[>=Latex, font=\small,
        block/.style={draw=utnblue!60, line width=0.5pt, fill=utnblue!8,
                      minimum width=2.1cm, minimum height=1cm, rounded corners=2pt,
                      blur shadow={shadow blur steps=4, shadow blur extra rounding=0pt,
                                   shadow xshift=0.4pt, shadow yshift=-0.6pt, shadow opacity=25}},
        signal/.style={->, thick, utnblue},
    ]
        \pgfplotsset{
            width=4cm, height=2.5cm, 
            axis lines=center,
            axis line style={->, thin},
            ticks=none,
            xmin=0, xmax=5, ymin=-1.5, ymax=1.5,
        }

        % =============================================
        % (a) AMPLIFICADOR
        % =============================================
        \node[font=\footnotesize\bfseries, anchor=west] at (-2.5, 0.75) {(a) Amplificador};
        \node[font=\footnotesize, utnblue] at (-0.5, 0.3) {$x(t)$};
        \node[font=\footnotesize, red!70!black] at (9, 0.3) {$y(t) = A\cdot x(t)$};

        % Entrada: onda pequeña
        \begin{axis}[at={(-0.5cm, -0.5cm)}, anchor=center, ymin=-0.6, ymax=0.6]
            \addplot[domain=0:5, samples=100, thick, utnblue] {0.3*sin(deg(2.5*x))};
        \end{axis}

        % Bloque
        \node[block] (amp) at (4.5, -0.5) {\textbf{Amplificador}};
        \node[font=\footnotesize] at (4.5, -1.3) {$\mathcal{T}\{x\} = A\cdot x$};

        % Flechas
        \draw[signal] (1.5, -0.5) -- (amp.west);

        % Salida: onda grande (misma forma, más alta)
        \begin{axis}[at={(9cm, -0.5cm)}, anchor=center]
            \addplot[domain=0:5, samples=100, thick, red!70!black] {1.2*sin(deg(2.5*x))};
        \end{axis}
        \draw[signal] (amp.east) -- (7, -0.5);

        % =============================================
        % (b) ECO
        % =============================================
        \begin{scope}[yshift=-2.6cm]
            \node[font=\footnotesize\bfseries, anchor=west] at (-2.5, 0.75) {(b) Sala con eco};
            \node[font=\footnotesize, utnblue] at (-0.5, 0.3) {$x(t)$};
            \node[font=\footnotesize, red!70!black] at (9, 0.3) {$y(t) = x(t) + 0.5\,x(t-\tau)$};

            % Entrada: pulso breve
            \begin{axis}[at={(-0.5cm, -0.5cm)}, anchor=center, ymin=-0.3, ymax=1.5]
                \addplot[domain=0:5, samples=200, thick, utnblue] 
                    {1.2*exp(-5*(x-1.5)^2)};
            \end{axis}

            % Bloque
            \node[block] (eco) at (4.5, -0.5) {\textbf{Sala con eco}};
            \node[font=\footnotesize] at (4.5, -1.3) {$\mathcal{T}\{x\} = x(t) + 0.5\,x(t-\tau)$};

            % Flechas
            \draw[signal] (1.5, -0.5) -- (eco.west);

            % Salida: pulso + eco retardado
            \begin{axis}[at={(9cm, -0.5cm)}, anchor=center, ymin=-0.3, ymax=1.5]
                \addplot[domain=0:5, samples=200, thick, red!70!black] 
                    {1.2*exp(-5*(x-1)^2) + 0.5*exp(-5*(x-3)^2)};
                % Etiquetas
                \node[font=\tiny, red!70!black] at (axis cs:1, 1.35) {original};
                \node[font=\tiny, red!70!black] at (axis cs:3.5, 0.65) {eco};
            \end{axis}
            \draw[signal] (eco.east) -- (7, -0.5);
        \end{scope}

        % =============================================
        % (c) SUSPENSIÓN
        % =============================================
        \begin{scope}[yshift=-5.2cm]
            \node[font=\footnotesize\bfseries, anchor=west] at (-2.5, 0.75) {(c) Suspensión de auto};
            \node[font=\footnotesize, utnblue] at (-0.5, 0.3) {Perfil del camino};
            \node[font=\footnotesize, red!70!black] at (9, 0.3) {Movimiento del chasis};

            % Entrada: perfil de camino irregular
            \begin{axis}[at={(-0.5cm, -0.5cm)}, anchor=center, ymin=-0.3, ymax=1.5]
                \addplot[domain=0:5, samples=300, thick, utnblue] 
                    {0.8*exp(-20*(x-1)^2) + 1.0*exp(-25*(x-2.5)^2) + 0.5*exp(-20*(x-3.8)^2)};
            \end{axis}

            % Bloque
            \node[block] (susp) at (4.5, -0.5) {\textbf{Suspensión}};
            \node[font=\footnotesize] at (4.5, -1.3) {$\mathcal{T}\{x\} = \;?$};

            % Flechas
            \draw[signal] (1.5, -0.5) -- (susp.west);

            % Salida: movimiento suavizado
            \begin{axis}[at={(9cm, -0.5cm)}, anchor=center, ymin=-0.3, ymax=1.5]
                \addplot[domain=0:5, samples=300, thick, red!70!black, smooth] 
                    {0.3*sin(deg(2*x))*exp(-0.3*x) + 0.25*exp(-2*(x-1.5)^2) + 0.2*exp(-2*(x-3)^2)};
            \end{axis}
            \draw[signal] (susp.east) -- (7, -0.5);
        \end{scope}

    \end{tikzpicture}
\end{document}
